% !TEX program = xelatex
% 高中数学培优专题：参数方程与极坐标
\documentclass[12pt,a4paper]{ctexart}

% ==================== 字体 ====================
\setCJKmainfont{Hiragino Sans GB}
\setCJKsansfont{Heiti SC}

% ==================== 宏包 ====================
\usepackage[top=2.5cm,bottom=2.5cm,left=2.5cm,right=2.5cm]{geometry}
\usepackage{amsmath,amssymb,amsthm}
\usepackage{tcolorbox}
\tcbuselibrary{skins,breakable}
\usepackage{enumitem}
\usepackage{fancyhdr}
\usepackage{graphicx}
\usepackage{xcolor}
\usepackage{tikz}
\usepackage{pgfplots}
\pgfplotsset{compat=1.18}
\usepackage{hyperref}
\hypersetup{colorlinks=true,linkcolor=blue!60!black}

% ==================== 颜色定义 ====================
\definecolor{mainblue}{RGB}{0,90,160}
\definecolor{darkblue}{RGB}{0,60,120}
\definecolor{lightblue}{RGB}{220,235,250}
\definecolor{examgreen}{RGB}{0,120,80}
\definecolor{lightgreen}{RGB}{220,245,230}
\definecolor{deforange}{RGB}{200,100,0}
\definecolor{lightorange}{RGB}{255,240,220}
\definecolor{warnred}{RGB}{180,30,30}
\definecolor{lightred}{RGB}{255,230,230}
\definecolor{practicepurple}{RGB}{100,50,150}
\definecolor{lightpurple}{RGB}{235,225,250}

% ==================== 页眉页脚 ====================
\pagestyle{fancy}
\fancyhf{}
\fancyhead[L]{\small\color{mainblue}\leftmark}
\fancyhead[R]{\small\color{mainblue}高中数学培优 · 参数方程与极坐标}
\fancyfoot[C]{\small\color{mainblue}\thepage}
\renewcommand{\headrulewidth}{0.4pt}
\renewcommand{\headrule}{\hbox to\headwidth{\color{mainblue}\leaders\hrule height \headrulewidth\hfill}}

% ==================== tcolorbox 样式 ====================
% 知识点框
\newtcolorbox{knowledge}[1][]{
  colback=lightblue,
  colframe=mainblue,
  fonttitle=\bfseries\large,
  title={#1},
  breakable,
  arc=3mm,
  boxrule=1pt,
  left=6pt,right=6pt,top=6pt,bottom=6pt,
}

% 定义框
\newtcolorbox{definition}[1][]{
  colback=lightorange,
  colframe=deforange,
  fonttitle=\bfseries,
  title={#1},
  breakable,
  arc=2mm,
  boxrule=1pt,
}

% 例题框
\newtcolorbox{example}[1][]{
  colback=white,
  colframe=examgreen,
  fonttitle=\bfseries,
  title={#1},
  breakable,
  arc=2mm,
  boxrule=1.2pt,
  left=8pt,right=8pt,top=6pt,bottom=6pt,
}

% 应用题框
\newtcolorbox{application}[1][]{
  colback=lightpurple,
  colframe=practicepurple,
  fonttitle=\bfseries,
  title={#1},
  breakable,
  arc=2mm,
  boxrule=1.2pt,
}

% 注意框
\newtcolorbox{note}{
  colback=lightred,
  colframe=warnred,
  fonttitle=\bfseries,
  title={\textbf{注意}},
  breakable,
  arc=2mm,
  boxrule=1pt,
}

% 练习框
\newtcolorbox{practice}[1][]{
  colback=lightgreen,
  colframe=examgreen,
  fonttitle=\bfseries,
  title={#1},
  breakable,
  arc=2mm,
  boxrule=1pt,
}

% ==================== 定理环境 ====================
\theoremstyle{definition}
\newtheorem*{remark}{备注}

% ==================== 正文 ====================
\begin{document}

\begin{center}
  {\LARGE\bfseries\color{darkblue} 高中数学培优专题}\\[6pt]
  {\Huge\bfseries\color{mainblue} 参数方程与极坐标}\\[10pt]
  {\large\color{gray} 参数方程的概念与应用\quad$\cdot$\quad 极坐标系及其应用\quad$\cdot$\quad 坐标变换}
\end{center}

\vspace{4pt}\hrule\vspace{12pt}

\tableofcontents
\newpage

% ============================================================
%                   第一部分：知识点梳理
% ============================================================
\section{知识点梳理}

% ---------- 1.1 参数方程的概念 ----------
\subsection{参数方程的基本概念}

\begin{definition}[\textsf{参数方程的定义}]
  在平面直角坐标系中，如果曲线 $C$ 上任意一点的坐标 $x,\,y$ 都可以表示为某个变量 $t$（称为\textbf{参数}）的函数
  \[
    \begin{cases}
      x = f(t), \\[4pt]
      y = g(t),
    \end{cases}
    \quad t \in D,
  \]
  则称上述方程组为曲线 $C$ 的\textbf{参数方程}，$t$ 称为\textbf{参数}。消去参数 $t$ 后得到的关于 $x,\,y$ 的方程 $F(x,y)=0$ 称为曲线 $C$ 的\textbf{普通方程}。
\end{definition}

\begin{note}
  参数方程与普通方程的互化时，必须注意\textbf{参数的取值范围}对 $x,\,y$ 取值范围的限制，两种方程表示的曲线可能不完全相同。
\end{note}

\begin{knowledge}[\textsf{常见曲线的参数方程}]
  \begin{enumerate}[leftmargin=2em]
    \item \textbf{直线}：过点 $P_0(x_0,y_0)$，倾斜角为 $\alpha$ 的直线参数方程
    \[
      \begin{cases}
        x = x_0 + t\cos\alpha, \\[3pt]
        y = y_0 + t\sin\alpha,
      \end{cases}
      \quad t \in \mathbb{R}.
    \]
    其中参数 $t$ 的几何意义：$|t|=|P_0P|$，当 $P$ 在 $P_0$ 上方（沿倾斜角方向）时 $t>0$，反之 $t<0$。

    \item \textbf{圆}：以 $C(a,b)$ 为圆心、$r$ 为半径的圆的参数方程
    \[
      \begin{cases}
        x = a + r\cos\theta, \\[3pt]
        y = b + r\sin\theta,
      \end{cases}
      \quad \theta \in [0,2\pi).
    \]

    \item \textbf{椭圆}：$\dfrac{x^2}{a^2}+\dfrac{y^2}{b^2}=1\;(a>b>0)$ 的参数方程
    \[
      \begin{cases}
        x = a\cos\varphi, \\[3pt]
        y = b\sin\varphi,
      \end{cases}
      \quad \varphi \in [0,2\pi).
    \]

    \item \textbf{双曲线}：$\dfrac{x^2}{a^2}-\dfrac{y^2}{b^2}=1$ 的参数方程
    \[
      \begin{cases}
        x = a\sec\varphi, \\[3pt]
        y = b\tan\varphi,
      \end{cases}
      \quad \varphi \neq \dfrac{\pi}{2}+k\pi,\;k\in\mathbb{Z}.
    \]

    \item \textbf{抛物线}：$y^2=2px$ 的参数方程
    \[
      \begin{cases}
        x = 2pt^2, \\[3pt]
        y = 2pt,
      \end{cases}
      \quad t \in \mathbb{R}.
    \]
  \end{enumerate}
\end{knowledge}

% ---------- 1.2 极坐标系 ----------
\subsection{极坐标系与极坐标方程}

\begin{definition}[\textsf{极坐标系}]
  在平面上取定点 $O$，称为\textbf{极点}；从 $O$ 出发的一条射线 $Ox$，称为\textbf{极轴}。对于平面上任意一点 $P$，用有序数对 $(\rho,\theta)$ 表示，其中 $\rho = |OP| \geqslant 0$ 称为\textbf{极径}，$\theta$ 为 $Ox$ 到 $OP$ 的角（\textbf{极角}），则 $(\rho,\theta)$ 称为点 $P$ 的\textbf{极坐标}。
\end{definition}

\begin{knowledge}[\textsf{极坐标与直角坐标的互化}]
  设极点与原点重合，极轴与 $x$ 轴正方向重合，则：
  \[
    \begin{cases}
      x = \rho\cos\theta, \\
      y = \rho\sin\theta,
    \end{cases}
    \qquad
    \begin{cases}
      \rho^2 = x^2 + y^2, \\[3pt]
      \tan\theta = \dfrac{y}{x}\;(x \neq 0).
    \end{cases}
  \]
\end{knowledge}

\begin{knowledge}[\textsf{常见曲线的极坐标方程}]
  \begin{enumerate}[leftmargin=2em]
    \item \textbf{过极点的直线}：$\theta = \alpha$（$\alpha$ 为常数）。

    \item \textbf{过点 $(a,0)$ 且垂直于极轴的直线}：$\rho\cos\theta = a$。

    \item \textbf{过点 $(a,\dfrac{\pi}{2})$ 且平行于极轴的直线}：$\rho\sin\theta = a$。

    \item \textbf{以极点为圆心、$r$ 为半径的圆}：$\rho = r$。

    \item \textbf{过极点、圆心在极轴上、半径为 $a$ 的圆}：$\rho = 2a\cos\theta$。

    \item \textbf{过极点、圆心在 $\theta=\dfrac{\pi}{2}$ 上、半径为 $a$ 的圆}：$\rho = 2a\sin\theta$。

    \item \textbf{圆锥曲线的统一极坐标方程}（焦点在极点，准线在左侧）：
    \[
      \rho = \frac{ep}{1 - e\cos\theta},
    \]
    其中 $e$ 为离心率，$p$ 为焦点到准线的距离。$0<e<1$ 为椭圆，$e=1$ 为抛物线，$e>1$ 为双曲线。
  \end{enumerate}
\end{knowledge}

\begin{note}
  在极坐标方程中，$\rho$ 一般取非负值，但在某些方程中允许 $\rho<0$，此时点 $(\rho,\theta)$ 与 $(-\rho,\theta+\pi)$ 表示同一点。
\end{note}

\newpage

% ============================================================
%                   第二部分：例题精讲
% ============================================================
\section{例题精讲}

% ---------- 例题1 ----------
\begin{example}[\textsf{例1\quad 参数方程与普通方程互化}]
  将下列参数方程化为普通方程，并说明曲线类型：
  \[
    \text{(1)}\;\begin{cases} x = 2\cos\theta, \\ y = 2\sin\theta, \end{cases}\quad \theta\in[0,2\pi);
    \qquad
    \text{(2)}\;\begin{cases} x = 1 + 2t, \\ y = -1 + t, \end{cases}\quad t\in\mathbb{R}.
  \]
\end{example}

\textbf{解}：

(1) 由 $x=2\cos\theta$，$y=2\sin\theta$，利用 $\cos^2\theta+\sin^2\theta=1$，得
\[
  \frac{x^2}{4} + \frac{y^2}{4} = 1 \implies x^2 + y^2 = 4.
\]
此曲线为以原点为圆心、$2$ 为半径的\textbf{圆}。

(2) 由 $x=1+2t$ 得 $t = \dfrac{x-1}{2}$，代入 $y=-1+t$：
\[
  y = -1 + \frac{x-1}{2} = \frac{x-3}{2} \implies x - 2y - 3 = 0.
\]
此曲线为一条\textbf{直线}。

% ---------- 例题2 ----------
\begin{example}[\textsf{例2\quad 直线参数方程的应用}]
  过点 $P_0(1,2)$ 且倾斜角为 $\dfrac{\pi}{3}$ 的直线 $l$ 与圆 $x^2+y^2=9$ 交于 $A$、$B$ 两点，求 $|PA|\cdot|PB|$。
\end{example}

\textbf{解}：

直线 $l$ 的参数方程为
\[
  \begin{cases}
    x = 1 + \dfrac{1}{2}t, \\[6pt]
    y = 2 + \dfrac{\sqrt{3}}{2}t,
  \end{cases}
  \quad t \in \mathbb{R}.
\]

代入圆的方程 $x^2+y^2=9$：
\[
  \left(1+\frac{t}{2}\right)^2 + \left(2+\frac{\sqrt{3}}{2}t\right)^2 = 9.
\]

展开整理：
\[
  1 + t + \frac{t^2}{4} + 4 + 2\sqrt{3}\,t + \frac{3t^2}{4} = 9,
\]
\[
  t^2 + (1+2\sqrt{3})\,t + 5 = 9,
\]
\[
  t^2 + (1+2\sqrt{3})\,t - 4 = 0.
\]

设 $t_1,\,t_2$ 为方程的两根，由韦达定理 $t_1 t_2 = -4$。由于参数 $t$ 的几何意义 $|P_0A|=|t_1|$，$|P_0B|=|t_2|$，所以
\[
  |PA|\cdot|PB| = |t_1| \cdot |t_2| = |t_1 t_2| = |-4| = \boxed{4}.
\]

% ---------- 例题3 ----------
\begin{example}[\textsf{例3\quad 圆的参数方程与最值}]
  已知 $P(x,y)$ 是圆 $C:(x-1)^2+(y-2)^2=4$ 上的动点，求 $x+y$ 的最大值和最小值。
\end{example}

\textbf{解}：

圆 $C$ 的参数方程为
\[
  \begin{cases}
    x = 1 + 2\cos\theta, \\
    y = 2 + 2\sin\theta,
  \end{cases}
  \quad \theta \in [0,2\pi).
\]

则
\[
  x + y = (1+2\cos\theta) + (2+2\sin\theta) = 3 + 2(\sin\theta + \cos\theta).
\]

利用辅助角公式：
\[
  \sin\theta + \cos\theta = \sqrt{2}\sin\!\left(\theta+\frac{\pi}{4}\right),
\]

所以
\[
  x+y = 3 + 2\sqrt{2}\sin\!\left(\theta+\frac{\pi}{4}\right).
\]

当 $\sin\!\left(\theta+\dfrac{\pi}{4}\right)=1$ 时，$(x+y)_{\max} = 3 + 2\sqrt{2}$；

当 $\sin\!\left(\theta+\dfrac{\pi}{4}\right)=-1$ 时，$(x+y)_{\min} = 3 - 2\sqrt{2}$。

\[
  \boxed{(x+y)_{\max} = 3+2\sqrt{2},\quad (x+y)_{\min} = 3-2\sqrt{2}.}
\]

% ---------- 例题4 ----------
\begin{example}[\textsf{例4\quad 椭圆参数方程的综合应用}]
  设椭圆 $\dfrac{x^2}{4}+y^2=1$ 上的点 $P$ 对应的参数为 $\varphi$，求点 $P$ 到直线 $x+y-4=0$ 的最小距离。
\end{example}

\textbf{解}：

椭圆的参数方程为 $P(2\cos\varphi,\,\sin\varphi)$。点 $P$ 到直线 $x+y-4=0$ 的距离为
\[
  d = \frac{|2\cos\varphi + \sin\varphi - 4|}{\sqrt{1^2+1^2}} = \frac{|2\cos\varphi+\sin\varphi-4|}{\sqrt{2}}.
\]

令 $f(\varphi) = 2\cos\varphi + \sin\varphi$，利用辅助角公式：
\[
  f(\varphi) = \sqrt{2^2+1^2}\sin(\varphi+\alpha) = \sqrt{5}\sin(\varphi+\alpha),
\]
其中 $\tan\alpha = 2$。故 $f(\varphi) \in [-\sqrt{5},\,\sqrt{5}]$。

因为 $\sqrt{5} \approx 2.236 < 4$，所以 $2\cos\varphi+\sin\varphi-4 < 0$ 恒成立，从而
\[
  d = \frac{4-(2\cos\varphi+\sin\varphi)}{\sqrt{2}}.
\]

$d$ 最小当 $2\cos\varphi+\sin\varphi$ 最大，即 $f(\varphi)=\sqrt{5}$ 时：
\[
  d_{\min} = \frac{4-\sqrt{5}}{\sqrt{2}} = \frac{(4-\sqrt{5})\sqrt{2}}{2} = \boxed{\frac{4\sqrt{2}-\sqrt{10}}{2}}.
\]

% ---------- 例题5 ----------
\begin{example}[\textsf{例5\quad 极坐标与直角坐标互化}]
  将极坐标方程 $\rho = 4\cos\theta$ 化为直角坐标方程，并说明曲线形状。
\end{example}

\textbf{解}：

方程 $\rho = 4\cos\theta$ 两边同乘 $\rho$：
\[
  \rho^2 = 4\rho\cos\theta.
\]

利用 $\rho^2=x^2+y^2$，$\rho\cos\theta = x$，得
\[
  x^2 + y^2 = 4x.
\]

配方：
\[
  (x-2)^2 + y^2 = 4.
\]

此曲线为以 $(2,0)$ 为圆心、$2$ 为半径的\textbf{圆}。

% ---------- 例题6 ----------
\begin{example}[\textsf{例6\quad 直线的极坐标方程}]
  在极坐标系中，求过点 $A\!\left(2,\dfrac{\pi}{3}\right)$ 且垂直于极轴的直线的极坐标方程。
\end{example}

\textbf{解}：

设直线上任意一点 $P(\rho,\theta)$。过 $A$ 作极轴的垂线，设垂足为 $H$。

在 $\triangle OPH$ 中（或直接利用投影关系），$P$ 到极轴的垂线段与 $A$ 到极轴的垂线段相等（因为直线垂直于极轴，直线上所有点的 $x$ 坐标相同）。

点 $A$ 的直角坐标为
\[
  x_A = 2\cos\frac{\pi}{3} = 1,\qquad y_A = 2\sin\frac{\pi}{3} = \sqrt{3}.
\]

所以直线的直角坐标方程为 $x=1$，即
\[
  \rho\cos\theta = 1 \implies \boxed{\rho = \frac{1}{\cos\theta} = \sec\theta.}
\]

% ---------- 例题7 ----------
\begin{example}[\textsf{例7\quad 圆锥曲线的统一极坐标方程}]
  已知椭圆的极坐标方程为 $\rho = \dfrac{5}{3-2\cos\theta}$，求该椭圆的离心率、长半轴和短半轴。
\end{example}

\textbf{解}：

将方程化为标准形式：
\[
  \rho = \frac{5}{3-2\cos\theta} = \frac{\frac{5}{3}}{1-\frac{2}{3}\cos\theta}.
\]

对比圆锥曲线统一方程 $\rho = \dfrac{ep}{1-e\cos\theta}$，得
\[
  e = \frac{2}{3}, \qquad ep = \frac{5}{3} \implies p = \frac{5}{3}\div\frac{2}{3} = \frac{5}{2}.
\]

离心率 $e = \dfrac{2}{3}$。

对于椭圆，$p = \dfrac{b^2}{c}$，又 $e = \dfrac{c}{a} = \dfrac{2}{3}$，即 $c = \dfrac{2a}{3}$。

由 $p = \dfrac{b^2}{c} = \dfrac{a^2-c^2}{c} = \dfrac{a^2 - \frac{4a^2}{9}}{\frac{2a}{3}} = \dfrac{\frac{5a^2}{9}}{\frac{2a}{3}} = \dfrac{5a}{6} = \dfrac{5}{2}$，

解得 $a = 3$。从而 $c = 2$，$b^2 = a^2-c^2 = 9-4 = 5$，$b = \sqrt{5}$。

\[
  \boxed{e = \frac{2}{3},\quad a = 3,\quad b = \sqrt{5}.}
\]

% ---------- 例题8 ----------
\begin{example}[\textsf{例8\quad 参数方程与韦达定理综合}]
  已知直线 $l$ 的参数方程为 $\begin{cases} x = 2+t, \\ y = \sqrt{3}\,t, \end{cases}$ $t\in\mathbb{R}$，曲线 $C$ 的极坐标方程为 $\rho=2\cos\theta$。设 $l$ 与 $C$ 交于 $A$、$B$ 两点，求 $|AB|$。
\end{example}

\textbf{解}：

\textbf{第一步}：将曲线 $C$ 化为直角坐标方程。$\rho=2\cos\theta$ 两边乘 $\rho$：
\[
  \rho^2 = 2\rho\cos\theta \implies x^2+y^2 = 2x \implies (x-1)^2+y^2=1.
\]

\textbf{第二步}：将直线参数方程代入曲线方程：
\[
  (2+t-1)^2 + (\sqrt{3}\,t)^2 = 1,
\]
\[
  (1+t)^2 + 3t^2 = 1,
\]
\[
  1+2t+t^2+3t^2 = 1,
\]
\[
  4t^2 + 2t = 0,
\]
\[
  2t(2t+1) = 0.
\]

解得 $t_1 = 0$，$t_2 = -\dfrac{1}{2}$。

\textbf{第三步}：此处参数方程的标准形式为 $x=x_0+t\cos\alpha$，$y=y_0+t\sin\alpha$。由 $x=2+t$，$y=\sqrt{3}\,t$ 可知，$\cos\alpha=1$，$\sin\alpha=\sqrt{3}$ 并非标准形式（不满足 $\cos^2\alpha+\sin^2\alpha=1$）。因此这里的 $t$ 不是标准参数，$|AB| \neq |t_1-t_2|$。

对应点坐标：$t_1=0$ 时 $A(2,0)$；$t_2=-\dfrac{1}{2}$ 时 $B\!\left(\dfrac{3}{2},\,-\dfrac{\sqrt{3}}{2}\right)$。

\[
  |AB| = \sqrt{\left(2-\frac{3}{2}\right)^2 + \left(0+\frac{\sqrt{3}}{2}\right)^2} = \sqrt{\frac{1}{4}+\frac{3}{4}} = \boxed{1}.
\]

% ---------- 例题9 ----------
\begin{example}[\textsf{例9\quad 极坐标中的弦长问题}]
  在极坐标系中，已知圆 $C:\rho=2\cos\theta$，直线 $l:\rho\cos\theta=1$。求圆 $C$ 被直线 $l$ 截得的弦长。
\end{example}

\textbf{解}：

化为直角坐标：圆为 $(x-1)^2+y^2=1$，圆心 $O'(1,0)$，半径 $r=1$；直线为 $x=1$。

圆心 $O'$ 到直线 $x=1$ 的距离 $d = |1-1| = 0$。

直线过圆心！因此截得的弦即为圆的直径：
\[
  |AB| = 2r = \boxed{2}.
\]

% ---------- 例题10 ----------
\begin{example}[\textsf{例10\quad 抛物线的参数方程应用}]
  已知抛物线 $y^2=4x$ 上一点 $P$ 的参数为 $t_0=2$（参数方程取 $\begin{cases} x=t^2,\\ y=2t \end{cases}$），求点 $P$ 处的切线方程。
\end{example}

\textbf{解}：

当 $t_0=2$ 时，$P(x_0,y_0) = (4,\,4)$。

\textbf{方法一（导数法）}：由 $y^2=4x$ 两边对 $x$ 求导：$2y\dfrac{dy}{dx}=4$，即 $\dfrac{dy}{dx}=\dfrac{2}{y}$。

在 $P(4,4)$ 处，斜率 $k = \dfrac{2}{4} = \dfrac{1}{2}$。

切线方程：$y - 4 = \dfrac{1}{2}(x-4)$，化简得
\[
  \boxed{x - 2y + 4 = 0.}
\]

\textbf{方法二（切线公式）}：抛物线 $y^2=2px$（此处 $2p=4$，$p=2$）上点 $(x_0,y_0)$ 处的切线方程为 $y_0 y = p(x+x_0)$，即
\[
  4y = 2(x+4) \implies 2y = x+4 \implies x-2y+4=0. \quad\checkmark
\]

% ---------- 例题11 ----------
\begin{example}[\textsf{例11\quad 参数方程求轨迹}]
  过定点 $A(a,0)$（$a>0$）作圆 $x^2+y^2=a^2$ 的动弦 $AB$，求弦 $AB$ 中点 $M$ 的轨迹方程。
\end{example}

\textbf{解}：

设圆的参数方程为 $B(a\cos\theta,\,a\sin\theta)$，则 $AB$ 的中点 $M(x,y)$ 满足
\[
  x = \frac{a+a\cos\theta}{2} = \frac{a(1+\cos\theta)}{2}, \qquad y = \frac{0+a\sin\theta}{2} = \frac{a\sin\theta}{2}.
\]

消去参数 $\theta$。由
\[
  \cos\theta = \frac{2x}{a}-1, \qquad \sin\theta = \frac{2y}{a}.
\]

利用 $\cos^2\theta+\sin^2\theta=1$：
\[
  \left(\frac{2x}{a}-1\right)^2 + \left(\frac{2y}{a}\right)^2 = 1,
\]
\[
  \frac{(2x-a)^2}{a^2} + \frac{4y^2}{a^2} = 1,
\]
\[
  (2x-a)^2 + 4y^2 = a^2,
\]
\[
  4\!\left(x-\frac{a}{2}\right)^2 + 4y^2 = a^2,
\]
\[
  \left(x-\frac{a}{2}\right)^2 + y^2 = \frac{a^2}{4}.
\]

中点 $M$ 的轨迹为以 $\left(\dfrac{a}{2},\,0\right)$ 为圆心、$\dfrac{a}{2}$ 为半径的圆。注意 $B \neq A$（即 $\theta\neq 0$），故轨迹去掉点 $(a,0)$。

\[
  \boxed{\left(x-\frac{a}{2}\right)^2 + y^2 = \frac{a^2}{4},\quad (x,y)\neq(a,0).}
\]

% ---------- 例题12 ----------
\begin{example}[\textsf{例12\quad 极坐标中面积的计算}]
  在极坐标系中，曲线 $C:\rho=2\cos\theta$，射线 $\theta=\dfrac{\pi}{6}$ 与曲线 $C$ 交于点 $A$，射线 $\theta=\dfrac{\pi}{3}$ 与曲线 $C$ 交于点 $B$，求 $\triangle OAB$ 的面积（$O$ 为极点）。
\end{example}

\textbf{解}：

$A$ 的极径：$\rho_A = 2\cos\dfrac{\pi}{6} = 2\cdot\dfrac{\sqrt{3}}{2} = \sqrt{3}$。

$B$ 的极径：$\rho_B = 2\cos\dfrac{\pi}{3} = 2\cdot\dfrac{1}{2} = 1$。

$\triangle OAB$ 的面积（两边为 $\rho_A,\,\rho_B$，夹角 $\Delta\theta = \dfrac{\pi}{3}-\dfrac{\pi}{6}=\dfrac{\pi}{6}$）：
\[
  S = \frac{1}{2}\,\rho_A\,\rho_B\,\sin\Delta\theta = \frac{1}{2}\cdot\sqrt{3}\cdot 1\cdot\sin\frac{\pi}{6} = \frac{1}{2}\cdot\sqrt{3}\cdot\frac{1}{2} = \boxed{\frac{\sqrt{3}}{4}}.
\]

\newpage

% ============================================================
%                   第三部分：综合应用题
% ============================================================
\section{综合应用题}

% ---------- 应用题1 ----------
\begin{application}[\textsf{应用题1\quad 参数方程与面积最值}]
  已知椭圆 $\dfrac{x^2}{4}+\dfrac{y^2}{3}=1$，$A$ 为椭圆的左顶点，$P$ 为椭圆上的动点，求 $\triangle AOP$（$O$ 为原点）面积的最大值。
\end{application}

\textbf{解}：

椭圆参数方程：$P(2\cos\varphi,\,\sqrt{3}\sin\varphi)$。左顶点 $A(-2,0)$。

以 $AO$ 为底，$|AO|=2$，$P$ 到 $x$ 轴的距离（即 $|y_P|$）为高：
\[
  S = \frac{1}{2}\cdot|AO|\cdot|y_P| = \frac{1}{2}\cdot 2\cdot|\sqrt{3}\sin\varphi| = \sqrt{3}\,|\sin\varphi|.
\]

当 $|\sin\varphi|=1$，即 $\varphi=\dfrac{\pi}{2}$ 或 $\dfrac{3\pi}{2}$ 时：
\[
  S_{\max} = \boxed{\sqrt{3}}.
\]

% ---------- 应用题2 ----------
\begin{application}[\textsf{应用题2\quad 极坐标中的焦点弦}]
  已知抛物线的极坐标方程为 $\rho = \dfrac{2}{1-\cos\theta}$，过焦点的弦 $AB$ 的极角分别为 $\theta_1=\dfrac{\pi}{3}$ 和 $\theta_2=\dfrac{4\pi}{3}$，求 $|AB|$。
\end{application}

\textbf{解}：

由 $\rho = \dfrac{2}{1-\cos\theta}$，对比 $\rho = \dfrac{ep}{1-e\cos\theta}$，得 $e=1$（抛物线），$ep=2$，$p=2$。

\[
  \rho_A = \frac{2}{1-\cos\frac{\pi}{3}} = \frac{2}{1-\frac{1}{2}} = 4,
\]
\[
  \rho_B = \frac{2}{1-\cos\frac{4\pi}{3}} = \frac{2}{1-\left(-\frac{1}{2}\right)} = \frac{2}{\frac{3}{2}} = \frac{4}{3}.
\]

焦点弦长：
\[
  |AB| = \rho_A + \rho_B = 4 + \frac{4}{3} = \boxed{\frac{16}{3}}.
\]

% ---------- 应用题3 ----------
\begin{application}[\textsf{应用题3\quad 参数方程与不等式}]
  已知 $P(x,y)$ 在椭圆 $\dfrac{x^2}{16}+\dfrac{y^2}{9}=1$ 上运动，求 $f(x,y)=3x+4y$ 的取值范围。
\end{application}

\textbf{解}：

令 $x = 4\cos\varphi$，$y = 3\sin\varphi$，则
\[
  f = 3(4\cos\varphi) + 4(3\sin\varphi) = 12\cos\varphi + 12\sin\varphi = 12(\sin\varphi+\cos\varphi).
\]

由辅助角公式 $\sin\varphi+\cos\varphi = \sqrt{2}\sin\!\left(\varphi+\dfrac{\pi}{4}\right)$：
\[
  f = 12\sqrt{2}\sin\!\left(\varphi+\frac{\pi}{4}\right).
\]

因为 $\sin\!\left(\varphi+\dfrac{\pi}{4}\right)\in[-1,1]$，所以
\[
  \boxed{f(x,y) \in [-12\sqrt{2},\;12\sqrt{2}].}
\]

% ---------- 应用题4 ----------
\begin{application}[\textsf{应用题4\quad 综合极坐标与参数方程}]
  在直角坐标系中，曲线 $C_1$ 的参数方程为 $\begin{cases} x=\cos\theta, \\ y=\sin\theta, \end{cases}$ 曲线 $C_2$ 的极坐标方程为 $\rho\cos\!\left(\theta-\dfrac{\pi}{4}\right)=\sqrt{2}$。求 $C_1$ 与 $C_2$ 的交点坐标。
\end{application}

\textbf{解}：

$C_1$ 为单位圆 $x^2+y^2=1$。

对 $C_2$，展开极坐标方程：
\[
  \rho\left(\cos\theta\cos\frac{\pi}{4}+\sin\theta\sin\frac{\pi}{4}\right) = \sqrt{2},
\]
\[
  \rho\left(\frac{\sqrt{2}}{2}\cos\theta + \frac{\sqrt{2}}{2}\sin\theta\right) = \sqrt{2},
\]
\[
  \frac{\sqrt{2}}{2}(\rho\cos\theta + \rho\sin\theta) = \sqrt{2},
\]
\[
  x + y = 2.
\]

联立 $\begin{cases} x^2+y^2=1, \\ x+y=2. \end{cases}$

由 $y=2-x$ 代入：$x^2+(2-x)^2=1$，$2x^2-4x+4=1$，$2x^2-4x+3=0$。

判别式 $\Delta = 16-24 = -8 < 0$，无实数解。

因此曲线 $C_1$ 与 $C_2$ \textbf{无交点}。

（直线 $x+y=2$ 与原点距离 $\dfrac{2}{\sqrt{2}}=\sqrt{2}>1$，不与单位圆相交，验证正确。）

\newpage

% ============================================================
%                   第四部分：练习题
% ============================================================
\section{练习题}

\begin{practice}[\textsf{练习1}]
  将参数方程 $\begin{cases} x = 3\cos\theta, \\ y = 5\sin\theta \end{cases}$ 化为普通方程。
\end{practice}

\begin{practice}[\textsf{练习2}]
  已知直线 $l$ 的参数方程为 $\begin{cases} x = 1 + t, \\ y = 2 - t \end{cases}$（$t\in\mathbb{R}$），求直线 $l$ 的普通方程及倾斜角。
\end{practice}

\begin{practice}[\textsf{练习3}]
  在极坐标系中，曲线 $\rho = 4\sin\theta$ 表示什么图形？求其直角坐标方程及圆心坐标。
\end{practice}

\begin{practice}[\textsf{练习4}]
  椭圆 $\dfrac{x^2}{9}+\dfrac{y^2}{4}=1$ 上的点 $P$ 对应参数 $\varphi$，求 $x+y$ 的取值范围。
\end{practice}

\begin{practice}[\textsf{练习5}]
  已知直线 $l$ 过点 $P_0(1,1)$，倾斜角 $\alpha=\dfrac{\pi}{4}$，求 $l$ 被圆 $x^2+y^2=4$ 截得的弦长。
\end{practice}

\begin{practice}[\textsf{练习6}]
  将极坐标方程 $\rho = 6\cos\theta - 2\sin\theta$ 化为直角坐标方程，并求圆心和半径。
\end{practice}

\begin{practice}[\textsf{练习7}]
  在极坐标系中，直线 $\theta=\dfrac{\pi}{4}$（$\rho\in\mathbb{R}$）与圆 $\rho=2\cos\theta$ 的交点的极坐标。
\end{practice}

\begin{practice}[\textsf{练习8}]
  已知双曲线的参数方程为 $\begin{cases} x = 2\sec\varphi, \\ y = 3\tan\varphi \end{cases}$，求其普通方程。
\end{practice}

% ============================================================
%                   第五部分：练习答案
% ============================================================
\newpage
\section{练习答案}

\textbf{练习1}：利用 $\cos^2\theta+\sin^2\theta=1$，得 $\dfrac{x^2}{9}+\dfrac{y^2}{25}=1$。这是一个椭圆。

\medskip

\textbf{练习2}：由 $x=1+t$，$y=2-t$，消去 $t$：$y=2-(x-1)=3-x$，即 $x+y-3=0$。斜率 $k=-1$，倾斜角 $\alpha=\dfrac{3\pi}{4}$。

\medskip

\textbf{练习3}：$\rho=4\sin\theta$ 两边乘 $\rho$：$\rho^2=4\rho\sin\theta$，即 $x^2+y^2=4y$，配方得 $x^2+(y-2)^2=4$。这是以 $(0,2)$ 为圆心、$2$ 为半径的圆。

\medskip

\textbf{练习4}：令 $x=3\cos\varphi$，$y=2\sin\varphi$，则 $x+y=3\cos\varphi+2\sin\varphi=\sqrt{13}\sin(\varphi+\alpha)$，其中 $\tan\alpha=\dfrac{3}{2}$。取值范围为 $[-\sqrt{13},\,\sqrt{13}]$。

\medskip

\textbf{练习5}：直线参数方程 $\begin{cases} x=1+\frac{\sqrt{2}}{2}t, \\ y=1+\frac{\sqrt{2}}{2}t \end{cases}$ 代入 $x^2+y^2=4$：$\left(1+\frac{\sqrt{2}}{2}t\right)^2+\left(1+\frac{\sqrt{2}}{2}t\right)^2=4$，$2\left(1+\frac{\sqrt{2}}{2}t\right)^2=4$，$\left(1+\frac{\sqrt{2}}{2}t\right)^2=2$，$1+\frac{\sqrt{2}}{2}t=\pm\sqrt{2}$。解得 $t_1=\sqrt{2}$，$t_2=-3\sqrt{2}$。弦长 $|t_1-t_2|=4\sqrt{2}$。但验证：$\left(\frac{\sqrt{2}}{2}\right)^2+\left(\frac{\sqrt{2}}{2}\right)^2=1$，是标准参数方程。$|AB|=|t_1-t_2|=|\sqrt{2}-(-3\sqrt{2})|=4\sqrt{2}$。

\medskip

\textbf{练习6}：$\rho=6\cos\theta-2\sin\theta$ 两边乘 $\rho$：$\rho^2=6\rho\cos\theta-2\rho\sin\theta$，即 $x^2+y^2=6x-2y$，配方得 $(x-3)^2+(y+1)^2=10$。圆心 $(3,-1)$，半径 $\sqrt{10}$。

\medskip

\textbf{练习7}：将 $\theta=\dfrac{\pi}{4}$ 代入 $\rho=2\cos\theta$：$\rho=2\cos\dfrac{\pi}{4}=\sqrt{2}$。交点为 $\left(\sqrt{2},\,\dfrac{\pi}{4}\right)$。注意 $\rho=2\cos\theta \geqslant 0$ 要求 $\theta\in\left[-\dfrac{\pi}{2},\dfrac{\pi}{2}\right]$，$\theta=\dfrac{\pi}{4}$ 满足条件，只有一个交点。

\medskip

\textbf{练习8}：由 $\sec\varphi=\dfrac{x}{2}$，$\tan\varphi=\dfrac{y}{3}$，利用 $\sec^2\varphi-\tan^2\varphi=1$：$\dfrac{x^2}{4}-\dfrac{y^2}{9}=1$。这是以原点为中心的双曲线。

\vfill

\begin{center}
  \color{gray}\rule{0.6\textwidth}{0.4pt}\\[6pt]
  {\small 高中数学培优专题\quad$\cdot$\quad 参数方程与极坐标\quad$\cdot$\quad 专题十二}
\end{center}

\end{document}
